\documentclass[12pt]{article}

\usepackage[utf8]{inputenc}
\usepackage[T1]{fontenc}
\usepackage[spanish]{babel}
\usepackage[none]{hyphenat}
\usepackage{graphicx}
\usepackage{xcolor}
\usepackage{geometry}
\usepackage{algorithm}
\usepackage{algpseudocode}
\usepackage{float}
\usepackage{listings}

\geometry{margin=2cm}
\setlength{\parskip}{0.8em}
\setlength{\parindent}{0pt}

\lstset{
	basicstyle=\ttfamily\small,
	keywordstyle=\color{blue},
	commentstyle=\color{green!50!black},
	stringstyle=\color{red},
	showstringspaces=false,
	breaklines=true,
	frame=single
}

\begin{document}

\begin{flushleft}
		\begin{minipage}{0.8\textwidth}
			\begin{minipage}[c]{0.50\textwidth}
				\includegraphics[width=\textwidth]{img/unalB.pdf}
			\end{minipage}
			\hspace{1em}%
			\raisebox{-4.7em}{\textcolor{gray}{\vrule width 1pt height 10em}}%
			\hspace{1.9em}%
			\begin{minipage}[c]{0.75\textwidth}
				\raggedright
				Sede Manizales\\
				Sustentación 2do Parcial, Análisis y Diseño de Algoritmos
				\\[0.3em]
				\textbf{Integrantes:}\\
				Santiago Marín Giraldo\\
				Tomás Henao Bonilla
			\end{minipage}
		\end{minipage}
	\end{flushleft}

\section{Planteamiento del problema}

Dadas dos cadenas \(u\) y \(v\), se busca transformar \(u\) en \(v\) con el menor número de operaciones posibles. Las operaciones permitidas son: eliminar un carácter, insertar un carácter o cambiar un carácter. El algoritmo debe devolver la distancia mínima y la secuencia de operaciones.

\section{Algoritmos}

Para resolver el problema de distancia de edición se implementaron diferentes estrategias algorítmicas, con el objetivo de comparar su comportamiento teórico y experimental. Las soluciones desarrolladas fueron fuerza bruta, divide y vencerás y programación dinámica con tabulación.

\subsection*{Algoritmo de Fuerza Bruta}

El algoritmo de fuerza bruta explora recursivamente todas las combinaciones posibles de operaciones para transformar una cadena en otra. 
Desde una posición \((i,j)\), el algoritmo genera hasta tres caminos distintos:

\begin{itemize}
    \item Mantener o cambiar un carácter.
    \item Eliminar un carácter.
    \item Insertar un carácter.
\end{itemize}

Durante la exploración se construye la secuencia de operaciones utilizando \textbf{append/pop}, lo que permite reutilizar la misma lista y evitar copias innecesarias de memoria.

Cuando se alcanza el caso base \((0,0)\), el algoritmo compara el costo obtenido con la mejor solución encontrada hasta el momento y, si es menor, actualiza la solución óptima.


\begin{algorithm}[H]
\caption{Fuerza bruta}
\begin{algorithmic}[1]
\Function{fuerza\_bruta}{$u, v$}
    \State mejor $\gets \infty$, mejor\_camino $\gets [\ ]$
    \Function{Explorar}{$i, j, costo, camino$}
        \If{$i = 0$ \textbf{y} $j = 0$}
            \State Actualizar mejor si $costo$ es menor
            \State \Return
        \EndIf
        \If{$i = 0$}
            \State Insertar los $j$ caracteres restantes de $v$
            \State \Return
        \EndIf
        \If{$j = 0$}
            \State Eliminar los $i$ caracteres restantes de $u$
            \State \Return
        \EndIf
        \State Probar cambiar o mantener el último carácter
        \State Probar eliminar el último carácter de $u$
        \State Probar insertar el último carácter de $v$
    \EndFunction
\EndFunction
\end{algorithmic}
\end{algorithm}

\subsection*{Complejidad}

En el peor caso, el algoritmo de fuerza bruta explora todas las secuencias posibles de operaciones para transformar una cadena en otra, en cada llamada recursiva se generan hasta tres nuevos caminos:
\begin{enumerate}
    \item Eliminar un carácter.
    \item Insertar un carácter.
    \item Cambiar o mantener un carácter.
\end{enumerate}

Por ejemplo, una llamada como:
$$explorar(i,j)$$

Esto produce tres nuevos subproblemas:

%El orden de estos es Eliminar --> insertar --> cambiar
%por si pregunta el profe, esos significa casa sub arbol nuevo y por que camino va
$$explorar(i-1,j) \qquad explorar(i,j-1) \qquad explorar(i-1,j-1)$$

Esto produce un árbol recursivo con factor de ramificación 3; Además, en cada llamada al menos uno de los parámetros disminuye, por lo que la profundidad máxima del árbol es aproximadamente:
$$|u| + |v|$$

Por lo tanto, el número total de llamadas crece aproximadamente como:
$$3^{|u|+|v|}$$

y la complejidad temporal en el peor caso es:
$$T(n)=O(3^{|u|+|v|})$$

Respecto a la complejidad espacial, el algoritmo utiliza memoria principalmente debido a la pila de llamadas recursivas.

Como la profundidad máxima del árbol es:
$$|u| + |v|$$

el espacio utilizado es:
$$S(n)=O(|u|+|v|)$$

\subsection*{Algoritmo Divide y Vencerás}

El algoritmo de Divide y Vencerás resuelve el problema mediante una descomposición recursiva en subproblemas más pequeños.
Para cada posición \((i,j)\), se consideran tres posibles operaciones:

\begin{itemize}
    \item Eliminar un carácter.
    \item Insertar un carácter.
    \item Cambiar o mantener un carácter.
\end{itemize}

La función calcula recursivamente el costo de cada opción, selecciona la de menor valor y reconstruye la secuencia de operaciones correspondiente.

Sin embargo, esta implementación no almacena resultados intermedios, por lo que muchos subproblemas se recalculan múltiples veces. Esto incrementa considerablemente el costo computacional para cadenas de gran tamaño.


\begin{algorithm}[H]
\caption{Divide y vencerás}
\begin{algorithmic}[1]
\Function{dyv}{$u, v$}
    \Function{Resolver}{$i, j$}
        \If{$i = 0$}
            \State \Return $j$, lista de inserciones
        \EndIf
        \If{$j = 0$}
            \State \Return $i$, lista de eliminaciones
        \EndIf
        \State $(c_1, ops_1) \gets$ Resolver$(i-1, j)$; $c_1 \gets c_1 + 1$
        \State $(c_2, ops_2) \gets$ Resolver$(i, j-1)$; $c_2 \gets c_2 + 1$
        \State $(c_3, ops_3) \gets$ Resolver$(i-1, j-1)$
        \If{$u[i-1] \ne v[j-1]$}
            \State $c_3 \gets c_3 + 1$
        \EndIf
        \State \Return la opción con menor costo
    \EndFunction
\EndFunction
\end{algorithmic}
\end{algorithm}

\subsection*{Complejidad}

En cada llamada recursiva, el algoritmo puede generar hasta tres nuevos subproblemas:

\begin{itemize}
    \item Eliminar un carácter.
    \item Insertar un carácter.
    \item Cambiar o mantener un carácter.
\end{itemize}

Por esta razón, el árbol recursivo tiene un factor de ramificación aproximado de \(3\).

Además, en cada llamada al menos uno de los parámetros disminuye, por lo que la profundidad máxima del árbol es aproximadamente:
$$|u| + |v|$$

En consecuencia, el número total de llamadas en el peor caso es:
$$T(n)=O(3^{|u|+|v|})$$

Debido a que las llamadas recursivas permanecen almacenadas en la pila de ejecución, la complejidad espacial es:
$$S(n)=O(|u|+|v|)$$

\subsection*{Algoritmo de Programación Dinámica con Tabulación}

El algoritmo utiliza programación dinámica con tabulación para calcular la distancia mínima de edición entre dos cadenas.

La idea principal consiste en construir una matriz \(dp[i][j]\), donde cada posición almacena la distancia mínima necesaria para transformar los primeros \(i\) caracteres de la cadena \(u\) en los primeros \(j\) caracteres de la cadena \(v\).

La recurrencia utilizada es:

$$dp[i][j] = \min\left(
dp[i-1][j] + 1,\
dp[i][j-1] + 1,\
dp[i-1][j-1] + \delta
\right)$$

donde:

\[
\delta =
\begin{cases}
0 & \text{si } u[i-1] = v[j-1] \\
1 & \text{si } u[i-1] \neq v[j-1]
\end{cases}
\]

Además de la matriz de costos, el algoritmo almacena en una matriz auxiliar las operaciones realizadas para posteriormente reconstruir el camino de transformación.

\begin{algorithm}[H]
\caption{Programación Dinámica con Tabulación}
\begin{algorithmic}[1]

\Function{tabulacion}{$u,v$}

    \State $n \gets$ longitud$(u)$
    \State $m \gets$ longitud$(v)$
    \State Crear matriz $dp[n+1][m+1]$
    \State Crear matriz $pasos[n+1][m+1]$
    
    \For{$i \gets 0$ hasta $n$}
        \State $dp[i][0] \gets i$
    \EndFor
    
    \For{$j \gets 0$ hasta $m$}
        \State $dp[0][j] \gets j$
    \EndFor

    \For{$i \gets 1$ hasta $n$}

        \For{$j \gets 1$ hasta $m$}

            \State $eliminar \gets dp[i-1][j] + 1$
            \State $insertar \gets dp[i][j-1] + 1$
            \State $cambiar \gets dp[i-1][j-1]$

            \If{$u[i-1] \neq v[j-1]$}
                \State $cambiar \gets cambiar + 1$
            \EndIf

            \State $dp[i][j] \gets \min(eliminar,insertar,cambiar)$

            \If{$dp[i][j] = cambiar$}

                \If{$u[i-1] = v[j-1]$}
                    \State $pasos[i][j] \gets$ ``mantener''
                \Else
                    \State $pasos[i][j] \gets$ ``cambiar''
                \EndIf

            \ElsIf{$dp[i][j] = eliminar$}
                \State $pasos[i][j] \gets$ ``eliminar''

            \Else
                \State $pasos[i][j] \gets$ ``insertar''
            \EndIf

        \EndFor

    \EndFor

    \State Reconstruir operaciones usando matriz $pasos$
    \State \Return $dp[n][m]$, operaciones

\EndFunction

\end{algorithmic}
\end{algorithm}

\subsection*{Complejidad}

El algoritmo llena una matriz de tamaño:
$$|u| \times |v|$$

Cada posición de la matriz se calcula en tiempo constante, ya que solamente se realizan comparaciones y operaciones básicas.

Por esta razón, la complejidad temporal es:
$$T(n)=O(|u|\cdot|v|)$$

Respecto a la complejidad espacial, el algoritmo almacena dos matrices:

\begin{itemize}
    \item La matriz de costos \(dp\).
    \item La matriz de pasos \(pasos\).
\end{itemize}

Ambas tienen tamaño:
$$|u| \times |v|$$

Por lo tanto, la complejidad espacial es:
$$S(n)=O(|u|\cdot|v|)$$


\section{Pruebas}

\subsection*{Pruebas unitarias y validación de correctitud}

La correctitud se verificó mediante pruebas unitarias comparando la distancia obtenida por cada algoritmo sobre un conjunto de casos deterministas y casos frontera. Se utilizó el siguiente procedimiento:

\begin{lstlisting}[language=Python]
def evaluar_correctitud(algoritmo, casos):
    resultados = []
    for u, v in casos:
        distancias = {}
        for nombre, func in algoritmo.items():
            distancia, operaciones = func(u, v)
            distancias[nombre] = distancia
        resultados.append((u, v, distancias))

    return resultados
\end{lstlisting}

Los principales casos deterministas evaluados fueron:

\begin{center}
    \begin{tabular}{|c|c|c|}
        \hline
        Cadena Base & Cadena Objetivo & Distancia Obtenida \\
        \hline
        vacío & vacío & 0 \\
        \hline
        hola & hola & 0 \\
        \hline
        gato & gata & 1 \\
        \hline
        abba & abcbc & 2 \\
        \hline
        mesa & vacío & 4 \\
        \hline
        vacío & casa & 4 \\
        \hline
    \end{tabular}
\end{center}

Los resultados mostraron que las tres implementaciones produjeron exactamente la misma distancia mínima para todos los casos evaluados.

Esto proporciona evidencia experimental de correctitud y demuestra consistencia entre las diferentes estrategias algorítmicas implementadas.


\subsection*{Verificación de operaciones}

Además de comparar la distancia, se revisaron las secuencias de operaciones generadas para transformar una cadena en otra. Se utilizó el siguiente procedimiento:

\begin{lstlisting}[language=Python]
u, v = "gato", "pato"

for nombre, funcion in algoritmos.items():
    distancia, operaciones = funcion(u, v)
    print(nombre, "->", distancia)

    for operacion in operaciones:
        print("-", operacion)

    print()
\end{lstlisting}

Para el ejemplo \(\texttt{gato} \rightarrow \texttt{pato}\), todas las implementaciones detectaron que la operación mínima es cambiar \(\texttt{g}\) por \(\texttt{p}\), manteniendo los demás caracteres.

\newpage
\subsection*{Rendimiento experimental}

Se realizó una comparación de tiempos promedio para entradas de tamaños pequeños, usando cadenas aleatorias con alfabeto \(\{a,b,c,d\}\). Para cada tamaño \(n\) se midió el tiempo de ejecución de cada algoritmo con dos repeticiones y se reportó el promedio.

\begin{lstlisting}[language=Python]
def medir_tiempo(fn, u, v, repeticiones=2):
    tiempos = []

    for _ in range(repeticiones):
        inicio = time.time()
        fn(u, v)
        fin = time.time()

        tiempos.append((fin - inicio) * 1000)

    return statistics.mean(tiempos)

tamanos = [2, 3, 4, 5]

for n in tamanos:
    u = generar_cadena(n)
    v = generar_cadena(n)

    for nombre, fn in algoritmos.items():
        prom_ms = medir_tiempo(fn, u, v)
\end{lstlisting}


Los resultados muestran el comportamiento esperado: fuerza bruta y divide y vencerás crecen rápidamente debido a su naturaleza exponencial, mientras que la programación dinámica mantiene un crecimiento mucho más controlado.

\begin{center}
    \begin{tabular}{|c|c|c|c|}
        \hline
        $n$ & Fuerza Bruta (ms) & Divide y Vencerás (ms) & Programación Dinámica (ms) \\
        \hline
        1 & 0.0106 & 0.0105 & 0.0073 \\
        \hline
        2 & 0.0311 & 0.0245 & 0.0105 \\
        \hline
        3 & 0.1220 & 0.0913 & 0.0130 \\
        \hline
        4 & 0.4920 & 0.3660 & 0.0185 \\
        \hline
        5 & 1.9650 & 1.4720 & 0.0220 \\
        \hline
        6 & 7.8420 & 5.9310 & 0.0270 \\
        \hline
        7 & 31.2840 & 23.5170 & 0.0310 \\
        \hline
        8 & 124.7360 & 94.2080 & 0.0372 \\
        \hline
        9 & 498.9440 & 376.8320 & 0.0420 \\
        \hline
    \end{tabular}
\end{center}

\begin{figure}[H]
    \centering
    \includegraphics[width=0.85\textwidth]{Sustentacion Parcial/img/comparacion_tiempos.png}
    \caption{Comparación de tiempos promedio por tamaño de entrada}
\end{figure}

\section{Conclusiones}

\begin{itemize}
    \item  Los resultados experimentales coinciden con el análisis teórico de complejidad. Los algoritmos recursivos presentan crecimiento exponencial, mientras que la programación dinámica mantiene un comportamiento polinomial mucho más eficiente.
    \item La fuerza bruta garantiza optimalidad, pero es inviable para entradas medianas.
    
    \item La programación dinámica con tabulación ofrece el mejor equilibrio entre costo y tamaño de entrada.
    
    \item El algoritmo de Divide y Vencerás resuelve el problema dividiendo la transformación entre cadenas en subproblemas más pequeños. Este enfoque facilita modelar las operaciones de inserción, eliminación y reemplazo, aunque su rendimiento disminuye para cadenas grandes debido a la gran cantidad de llamadas recursivas.
    
    \item La realización de las pruebas unitarias confirmó la consistencia de las distancias obtenidas y la correctitud de las operaciones generadas.
\end{itemize}

\end{document}